% MixSan: Enhancing Address-Based Memory Sanitizers with Pointer-Based Identity Verification
% Target: CMC - Computers, Materials & Continua (Tech Science Press)
\documentclass[cmc,article,submit,moreauthors,pdftex]{Definitions/tsp}

\input{Definitions/package}
\input{Definitions/unicode}

% For double-blind review, author info is redacted
\pubvolume{0}
\issuenum{0}
\elocationid{0}
\articlenumber{0}
\pubyear{2026}
\copyrightyear{2026}
\datereceived{Day Month Year}
\dateaccepted{Day Month Year}

\Title{MixSan: Enhancing Address-Based Memory Sanitizers with Pointer-Based Identity Verification}

\Author{Firstname Lastname\textsuperscript{1}}

\AuthorNames{Firstname Lastname}

\address{%
\textsuperscript{1}Affiliation 1 (Department/School/Faculty/Campus/\dots, University/\dots), City, Country
}%

\corres{Corresponding Author: Author's Name. Email: author@institute.xxx}

%===========================================================================
% Abstract (200-400 words, no formatting commands or paragraph breaks)
%===========================================================================
\abstract{
During software testing, memory errors in C/C++ often corrupt program state silently, leaving the program vulnerable to exploitation. Memory Sanitizer is designed to break this silence, transforming silent memory corruption into precise diagnostic signals at the point of occurrence which enable efficient bug localization, analysis, and repair. Since AddressSanitizer (ASan, 2012), address-based memory sanitizer, which detect violations by checking whether an accessed address is currently valid, have become one of the mainstream technical paradigms due to their practical performance and compatibility. However, address-based sanitizers share three common limitations. First, address-based spatial vulnerability detection judges legality solely by the target address, so redzone-skipping spatial vulnerabilities are undetectable due to lack of pointer semantics. Second, reuse-delay-based temporal vulnerability detection is unable to detect post-reuse temporal vulnerabilities. Third, the reuse-delay mechanism itself introduces substantial space and time overhead, particularly in multi-threaded scenario. In this paper, we propose a taxonomy, which serves as not only a classification framework but also a design methodology, based on the two key problems of memory sanitizer: validity encoding and violation detection. Based on the enlightening insights from the taxonomy, MixSan is developed, a prototype built upon state-of-the-art address-based sanitizer RangeSanitizer(RSan, 2025), which overcomes these three limitations by combining address-based detection with pointer-based identity verification, without introducing a reuse-delay mechanism. MixSan assigns each allocated object a 6-bit identity tag, stamping it into both the object's metadata and bits 62--57 of every pointer referencing it, which is transparent to address translation under Intel LAM U57. At each memory access, MixSan extracts the implicit size-class tag from the pointer to locate the object base, loads a single 8-byte metadata word storing the object end address, and evaluates the unified arithmetic check (Meta - (Ptr + n)) >> 43 != 0 in a branchless register sequence. This single formula simultaneously detects spatial overflows and temporal violations: an out-of-bounds access wraps the unsigned subtraction to produce a non-zero shift result; a dangling pointer---whose metadata was zeroed on free---or a pointer with corrupted MemTag bits triggers the same arithmetic effect. The lock-free allocator eliminates the quarantine entirely, using rdtsc()-based fast tag assignment and bitwise-NOT invalidation on free. On SPEC CPU2006, MixSan achieves 1.67x geomean runtime overhead, comparable to RSan's 1.69x, while using 43\% less memory (1.32x vs.\ 2.32x MaxRSS). Under 12-thread contention, the lock-free allocator delivers 712M ops/s versus RSan's serialized 5M ops/s. On a 102-test benchmark suite covering redzone-bypassing spatial errors and temporal-after-reuse errors, MixSan detects 98.0\% of bugs versus RSan's 2.0\%. These results demonstrate that pointer-based identity verification, when systematically integrated with address-based detection guided by design-semantic analysis, simultaneously improves detection capability, eliminates multi-thread bottlenecks, and maintains practical single-thread performance.
}

\keyword{memory safety; sanitizer; address-based detection; pointer-based identity; validity model; Intel LAM; MemTag; lock-free allocator}

%%%%%%%%%%%%%%%%%%%%%%%%%%%%%%%%%%%%%%%%%%%
\begin{document}

\section{Introduction} \label{sect:intro}

Memory safety errors in C and C++ programs---buffer overflows, use-after-free (UAF), double-free---have been the dominant root cause of exploitable software vulnerabilities for over three decades. Studies by Microsoft and Google consistently report that memory safety violations account for approximately 70\% of all critical security bugs in large-scale production codebases~\cite{msrc2019,chromium2020}. Sanitizers---compiler-instrumented runtime detectors that transform silent memory corruption into precise diagnostic signals at the point of occurrence---serve as the primary defense against such errors during dynamic program analysis. Among the diverse sanitizer designs proposed in the literature, \textbf{address-based} approaches---those that detect violations by checking whether the accessed address is currently legal, without knowledge of the accessing pointer's provenance---have achieved the widest practical adoption due to their favorable performance and relative implementation simplicity.

A memory safety sanitizer must address two fundamental design questions. The first concerns the \textbf{validity model}: how to design and manage the metadata that encodes the validity of pointers and memory addresses? Existing approaches employ either \emph{boundary metadata}---describing the spatial extent of an object (its base, size, or end address)---or \emph{identity-tag metadata}---encoding an identifier carried by both the object and any pointer that references it. The second concerns the \textbf{detection paradigm}: how to catch an access that violates validity constraints? Approaches are either \emph{pointer-based}---verifying at each access that the pointer's identity matches the target object's identity---or \emph{address-based}---checking whether the accessed address is currently legal, with no knowledge of the pointer's origin. The combinations of these two binary choices define a 2$\times$2 design space spanning four quadrants: Q1 (Boundary $\times$ Pointer-based), Q2 (Boundary $\times$ Address-based), Q3 (Identity-Tag $\times$ Pointer-based), and Q4 (Identity-Tag $\times$ Address-based). We develop this taxonomy not merely as a retrospective classification framework, but as a \emph{design methodology}: by revealing the comparative advantage of each quadrant---which safety properties each combination is inherently suited for, and which costs are intrinsic versus incidental---it guides the systematic construction of hybrid sanitizers.

\textbf{Address-based sanitizers: state of the art and its limitations.}
AddressSanitizer (ASan)~\cite{serebryany2012addresssanitizer}, the most widely deployed sanitizer and integrated into both GCC and LLVM, operates primarily in Q2: it places redzones (inaccessible memory regions) around each allocated object and instruments memory accesses with checks against per-byte shadow memory encoding address accessibility. For temporal safety, ASan adopts a Q4 quarantine---freed objects are held in a deferred-reuse queue so that dangling pointer dereferences land on inaccessible memory and trigger faults. RangeSanitizer (RSan)~\cite{gorter2025rsan}, published at USENIX Security 2025, refines this architecture: instead of per-byte shadow memory, RSan stores a per-object end-address as boundary metadata, enabling precise range comparisons that determine whether any access---at any offset---exceeds the object's spatial bound. For temporal safety, RSan employs a 256MB quarantine ring buffer protected by a global \texttt{pthread\_mutex\_lock}. On SPEC CPU2006, RSan achieves 1.69$\times$ geomean runtime overhead, making it the fastest sanitizer providing both spatial and temporal safety.

However, the architecture shared by ASan and RSan---boundary metadata for spatial detection, reuse-delay for temporal detection, and purely address-based checking for both---has three inherent limitations that stem directly from its design semantics:

\begin{enumerate}
\item \textbf{Redzone-bypassing spatial errors evade detection.} Address-based detection judges legality solely by the target address. An out-of-bounds access whose offset is large enough to skip the redzone and land within an adjacent valid object---a ``far-jump'' OOB---produces a legal target address. The detector, having no knowledge of which pointer initiated the access, cannot distinguish this from a legitimate access to the adjacent object. This is not a corner case: any sanitizer that relies purely on address-based checking is structurally blind to pointer-identity-violating spatial errors.

\item \textbf{Temporal safety expires after memory reuse.} The quarantine's detection guarantee ends when the freed memory slot is reallocated. Once a new object occupies the slot, a dangling pointer to the old object encounters valid metadata and passes the spatial check---the detector sees a legal address within a valid object. The detection window is probabilistic, bounded by quarantine capacity and program allocation rate. In contrast, a pointer carrying a persistent identity tag would detect the mismatch regardless of how many reuse cycles have occurred.

\item \textbf{Centralized quarantine serializes multi-threaded allocation.} RSan's 256MB ring buffer is protected by a single \texttt{pthread\_mutex\_lock}, serializing all free operations. Our measurements show allocator throughput collapsing from 730M to 5M ops/s under 12-thread contention---a cost invisible in single-threaded SPEC CPU2006 benchmarks. Any sanitizer relying on a centralized quarantine inherits this structural scalability ceiling.
\end{enumerate}

\textbf{The opportunity: pointer-based identity verification.}
The root cause of all three limitations is the absence of \emph{pointer-based identity verification} (Q3). If each pointer carried a persistent identity tag that could be verified against the target object at each access, then: (1) a far-jump OOB from a pointer tagged for object A would be caught when accessing object B, regardless of B's validity, because the tag mismatch reveals the pointer's origin; (2) temporal safety would be deterministic rather than window-based, since the tag mismatch persists through arbitrary reuse cycles; and (3) the quarantine could be eliminated entirely, since temporal detection no longer depends on delayed reuse. Our taxonomy reveals precisely this insight: Q3 (Identity-Tag $\times$ Pointer-based) is the systematically underexplored quadrant, complementary to Q2, whose adoption modern hardware features have newly made practical.

Intel Linear Address Masking (LAM)~\cite{intel2020lam} is the key enabler. LAM causes the processor to ignore bits 62--57 of a pointer during address translation (U57 mode), allowing these bits to carry arbitrary metadata without requiring explicit masking on every dereference. This provides a zero-overhead substrate for embedding identity tags in pointers---a capability that earlier tag-based sanitizers (e.g., CETS's key-lock mechanism~\cite{nagarakatte2010cets}) could not exploit because they required explicit tag extraction and comparison via disjoint metadata table lookups.

\textbf{MixSan: systematic integration of Q2 and Q3.}
We present MixSan, a memory sanitizer that systematically integrates pointer-based identity verification (Q3) into RSan's address-based spatial detection (Q2). MixSan embeds a 6-bit MemTag version number in pointer bits 62--57 via Intel LAM U57. On each \texttt{malloc()}, the allocator assigns the object's current tag to the returned pointer and stores the tag in the object's metadata. On each \texttt{free()}, the object's tag is incremented by +7---chosen because $\gcd(7,64)=1$, ensuring the tag cycles through all 63 non-reserved values before repeating, yielding a 1/63 collision probability per slot reuse. Every instrumented memory access passes through a \textbf{unified three-stage inline check}: (i) a SizeTag gate (Q4 fast path) filters obviously-in-bounds accesses using a single unsigned comparison; (ii) a MemTag comparison (Q3) verifies that the pointer's embedded tag matches the target object's current tag; (iii) a spatial bound comparison (Q2) confirms the access offset lies within the object. The three stages share a single metadata load and lower to a branchless x86 sequence. The \textbf{lock-free allocator} uses \texttt{rdtsc()} on the hot path (~20 cycles, register-only) for fast timestamp-based allocation and a bitwise-NOT tag invalidation (\texttt{*meta = \textasciitilde meta}) on free---avoiding the read-before-write dependency of zero-clearing. The quarantine ring buffer and its \texttt{pthread\_mutex\_lock} are eliminated entirely.

\textbf{Results.}
On SPEC CPU2006, MixSan achieves a geomean runtime overhead of 1.67$\times$, comparable to RSan's 1.69$\times$, despite performing strictly more work per check (tag comparison in addition to range verification). This parity is achieved through the optimized unified check algorithm. Memory consumption is 43\% lower than RSan (1.32$\times$ vs.\ 2.32$\times$ MaxRSS) due to the elimination of the 256MB quarantine buffer. Under 12-thread contention, the lock-free allocator delivers 712M ops/s versus RSan's serialized 5M ops/s---a $>$140$\times$ improvement. On a 102-test benchmark suite explicitly covering the detection blind spots of address-based sanitizers---redzone-bypassing spatial errors and temporal errors after memory reuse---MixSan detects 100/102 bugs (98.0\%) versus RSan's 2/102 (2.0\%). The two undetected cases are one intra-object overflow (a universal limitation of redzone-based sanitizers) and one MemTag collision at exactly 63 reuse cycles.

\textbf{Contributions.} This paper makes three contributions:

\begin{enumerate}
\item \textbf{A design-semantic taxonomy of memory sanitizers (Section~\ref{sect:background}).} We formalize a 2$\times$2 design space along two semantic axes---validity model (boundary-based vs.\ identity-tag-based metadata for encoding pointer and address validity) and detection paradigm (pointer-based identity verification vs.\ address-based legality checking). Unlike prior taxonomies that classify by implementation characteristics (metadata storage location, software/hardware substrate, or per-pointer/per-object binding), our taxonomy classifies by \emph{what the metadata means} and \emph{what the detection mechanism knows about the accessing pointer}---properties that directly determine detection capability and performance cost. We demonstrate that this taxonomy functions as a design methodology: it reveals that Q3 (Identity-Tag $\times$ Pointer-based) is the systematically underexplored quadrant that modern hardware features (LAM) have made newly practical, and that fusing Q2+Q3 systematically eliminates the limitations inherent to purely address-based architectures.

\item \textbf{MixSan: enhancing address-based detection with pointer-based identity verification (Sections~\ref{sect:design}--\ref{sect:implementation}).} MixSan systematically integrates Q3 pointer-based identity verification into RSan's Q2 address-based spatial detection. The 6-bit MemTag in LAM-masked pointer bits provides deterministic temporal safety without a quarantine; the lock-free allocator eliminates deferred-reuse overhead and scales near-linearly under multi-threaded contention; and the unified three-stage inline check (SizeTag gate $\rightarrow$ MemTag match $\rightarrow$ spatial bound) fuses both detection paradigms within a single branchless x86 sequence. On SPEC CPU2006, MixSan achieves 1.67$\times$ geomean overhead (vs.\ RSan's 1.69$\times$) with 43\% lower memory. The lock-free allocator delivers $>$140$\times$ higher throughput under 12-thread contention (712M vs.\ 5M ops/s).

\item \textbf{A benchmark suite for capability-boundary evaluation (Section~\ref{sec:eval:detection}).} We contribute a 102-test benchmark suite covering two detection scenarios absent from existing suites: temporal errors after memory reuse (UAF and double-free after slot reallocation, exercising the gap left by quarantine-based detectors) and redzone-bypassing spatial errors (OOB accesses that skip redzones and land within adjacent valid objects, exercising the gap left by purely address-based detectors). This suite enables principled comparison between deferred-reuse and versioned-identity sanitizer designs. MixSan detects 98.0\% of these bugs versus RSan's 2.0\%.
\end{enumerate}

%===========================================================================
% For double-blind review: all author info, funding, and identifiers are redacted
% Insert proper \acknowledgement, \funding, \authorcontributions,
% \availabilityofdataandmaterials, \ethicsapproval, \conflictsofinterest as needed
%===========================================================================

\end{document}
